%% Rewritten 2026-09-08 after the Marvin 26.1 regeneration and the reversibility
%% overhaul. Every number here is measured against the shipped files; see the
%% inline comments for the command that reproduces each. Detail in
%% Supplementary Methods S3.
\subsection{Multi-source thermodynamics}\label{sec:methods-thermo}

The 2020 release combined eQuilibrator with a group-contribution fallback. This
update ships four sources --- eQuilibrator~3.0~\cite{beber2022}, group
contribution rebuilt from MFAToolkit, dGPredictor~\cite{wang2021} retrained on
ModelSEED structures, and experimental values where they exist --- each stored
as its own \texttt{[energy, uncertainty, direction]} record so that no source is
overwritten by another. All values are reported at pH~7.0, ionic strength
0.25\,M, pMg~3.0 and 298.15\,K; Supplementary Methods~S3 explains why that
combination matches neither eQuilibrator preset. Only the experimental values
carry a published direction: the three predictors were released as energy
estimators, and every direction reported here is our inference from their
energies, not a claim by their authors.

\paragraph{Protonation.} eQuilibrator's Legendre transform requires a
\emph{macroscopic} pKa ladder per compound: an ordered list in which
consecutive species differ by one proton. The count of entries above the
reported pH sets the proton count of the major microspecies, so the composition
of the ladder changes the transformed energy directly, not merely its
precision.

%% Changed 2026-09-08. The previous build resolved ladders through a cascade
%% (alberty > iupac > cache > marvin > molgpka, ChemAxon tiers gated on MolGpKa
%% degeneracy) and reported 62% open by compound / 95% of reactions ChemAxon.
%% Those figures described a build that no longer ships. Current provenance is
%% derived, not transcribed, from
%% Biochemistry/Thermodynamics/ProtonationEvidence/pka_provenance.tsv.
Earlier builds assembled these ladders from whichever source could supply one,
preferring measured compilations, then an open predictor, and consulting a
licensed tool only where the open predictor provably degenerates. We have
abandoned that arrangement. Ladders from different models are not
interchangeable within one dataset: where a compound and its reaction partner
come from different models, a disagreement about how many sites lie above the
reported pH does not cancel across the reaction, and each unbalanced proton
displaces the transformed energy. Mixing sources to maximise coverage buys
coverage with an error invisible in any single compound's record.

Every ladder behind the shipped energies is therefore ChemAxon-derived, though
not all from one run: the builder reuses a pinned-release row wholesale wherever
the cache already holds our molecule, and several such compounds --- small
inorganics among them --- cannot be rebuilt at all, because group decomposition
rejects them. Results reports how far the regeneration actually reached. The
consequence is stated rather than elided: this layer is not open-source, and
alone in the database it cannot be regenerated without a licensed tool. We make
the alternative auditable instead, shipping the open predictions beside the ones
we used --- MolGpKa~\cite{pan2021} for every compound with a structure, and the
measured ladders of Alberty~\cite{alberty2003}, versioned per structure source
in an extended pKa dictionary. A reader can rebuild the layer on an open basis
and see where the models part company. We publish both; we do not mix them.

\paragraph{Uncertainty.} The sources report on scales that are not comparable
(Figure~\ref{fig:thermo}C): median 0.63\,kcal\,mol$^{-1}$ for eQuilibrator
against 10.41 for group contribution and 17.01 for dGPredictor. Against 797
experimentally-anchored reactions, group contribution overstates its error and
eQuilibrator understates it. eQuilibrator's figure is the
component-contribution covariance and carries no protonation term at all, so it
does not bound the effect described above. dGPredictor's, by contrast, is
well-calibrated: across 20,093 reactions with a tight eQuilibrator reference,
the residuals scaled by its own reported sigma have an RMS of 1.15, with 90\%
inside 1.69 sigma against the 1.645 a Gaussian predicts. Its error bars are
wide because it is uncertain, not because they are inflated.

\paragraph{Direction.} Group contribution keeps the rule set of Jankowski
\textit{et al.}~\cite{jankowski2008} unchanged from 2020, so that series stays
comparable across releases. eQuilibrator and dGPredictor share one rule set
built on the Noor reversibility index~\cite{noor2012}: a reaction is called
irreversible when $|\ln\Gamma|$ exceeds $\ln(1000)$ by more than one propagated
standard deviation~\cite{beber2022}.

Our gate is deliberately more conservative than the index as published. Noor
\textit{et al.} define $\Gamma$ as a point estimate tested against a threshold,
so every reaction falls on one side or the other; we require the propagated
uncertainty to clear the threshold as well, which withholds a directional call
from 58.9\% of dGPredictor's reactions and 15.7\% of eQuilibrator's that a
point-estimate test would decide. Two consequences follow. We report fewer
directions than the published index would on the same energies, and we
distinguish \emph{reversible} from \emph{undetermined}, a distinction the index
does not make. The exception is a reaction whose uncertainty is small relative
to the threshold and yet still straddles it: there the margin cannot resolve
which side the reaction falls on --- straddling means the estimate lies within
one standard deviation of the threshold by construction --- so we report it
reversible rather than undetermined, on the reading that a reaction poised at
the limit of physiological concentration control is one the cell can push either
way. Two guards precede it. A source that
declines to estimate is caught before any rule reads the number it published
alongside that refusal. And ATP synthase is assigned reversible structurally,
because the proton-motive force is its entire driving force and no source here
computes it --- without that rule the index returns $|\ln\Gamma| = 11.70 \pm
0.09$ and calls the pump irreversible in whichever direction it happens to be
transcribed. That caveat generalises: transport reactions are scored from
stoichiometry and energy alone, and neither the membrane potential nor the
transmembrane pH gradient enters any value in this release.

We also separate two claims the 2020 release conflated. A reaction is marked
reversible only where the evidence positively supports it --- where the error
bar fits inside the reversible band --- and undetermined where the uncertainty
straddles the threshold. Both previously shipped as reversible.
